\documentclass[10pt,a4paper]{amsart}
\usepackage{amsmath,amssymb,amsthm,graphicx}
\usepackage{verbatim}
\newenvironment{reference}{\comment}{\endcomment}
%\newenvironment{reference}{}{}
\newenvironment{slogan}{\comment}{\endcomment}
\newenvironment{history}{\comment}{\endcomment}
% For commutative diagrams we use Xy-pic
\usepackage[all]{xy}

% We use 2cell for 2-commutative diagrams.
\xyoption{2cell}
% Theorem environments.
%
\theoremstyle{plain}
\newtheorem{theorem}[subsection]{Theorem}
\newtheorem{proposition}[subsection]{Proposition}
\newtheorem{lemma}[subsection]{Lemma}

\theoremstyle{definition}
\newtheorem{definition}[subsection]{Definition}
\newtheorem{example}[subsection]{Example}
\newtheorem{exercise}[subsection]{Exercise}
\newtheorem{situation}[subsection]{Situation}

\theoremstyle{remark}
\newtheorem{remark}[subsection]{Remark}
\newtheorem{remarks}[subsection]{Remarks}

\numberwithin{equation}{subsection}

% Macros
%
\def\lim{\mathop{\mathrm{lim}}\nolimits}
\def\colim{\mathop{\mathrm{colim}}\nolimits}
\def\Spec{\mathop{\mathrm{Spec}}}
\def\Hom{\mathop{\mathrm{Hom}}\nolimits}
\def\Ext{\mathop{\mathrm{Ext}}\nolimits}
\def\SheafHom{\mathop{\mathcal{H}\!\mathit{om}}\nolimits}
\def\SheafExt{\mathop{\mathcal{E}\!\mathit{xt}}\nolimits}
\def\Sch{\mathit{Sch}}
\def\Mor{\mathop{\mathrm{Mor}}\nolimits}
\def\Ob{\mathop{\mathrm{Ob}}\nolimits}
\def\Sh{\mathop{\mathit{Sh}}\nolimits}
\def\NL{\mathop{N\!L}\nolimits}
\def\CH{\mathop{\mathrm{CH}}\nolimits}
\def\proetale{{pro\text{-}\acute{e}tale}}
\def\etale{{\acute{e}tale}}
\def\QCoh{\mathit{QCoh}}
\def\Ker{\mathop{\mathrm{Ker}}}
\def\Im{\mathop{\mathrm{Im}}}
\def\Coker{\mathop{\mathrm{Coker}}}
\def\Coim{\mathop{\mathrm{Coim}}}

% Boxtimes
%
\DeclareMathSymbol{\boxtimes}{\mathbin}{AMSa}{"02}

%
% Macros for moduli stacks/spaces
%
\def\QCohstack{\mathcal{QC}\!\mathit{oh}}
\def\Cohstack{\mathcal{C}\!\mathit{oh}}
\def\Spacesstack{\mathcal{S}\!\mathit{paces}}
\def\Quotfunctor{\mathrm{Quot}}
\def\Hilbfunctor{\mathrm{Hilb}}
\def\Curvesstack{\mathcal{C}\!\mathit{urves}}
\def\Polarizedstack{\mathcal{P}\!\mathit{olarized}}
\def\Complexesstack{\mathcal{C}\!\mathit{omplexes}}
% \Pic is the operator that assigns to X its picard group, usage \Pic(X)
% \Picardstack_{X/B} denotes the Picard stack of X over B
% \Picardfunctor_{X/B} denotes the Picard functor of X over B
\def\Pic{\mathop{\mathrm{Pic}}\nolimits}
\def\Picardstack{\mathcal{P}\!\mathit{ic}}
\def\Picardfunctor{\mathrm{Pic}}
\def\Deformationcategory{\mathcal{D}\!\mathit{ef}}
\usepackage{cleveref}
\crefname{originalreference}{original reference}{original references}
\Crefname{originalreference}{Original reference}{Original references}
\makeatletter
\newlabel{categories-section-posets-limits}{{categories-section-posets-limits}{source}{}{sourcepointer-1}{}}
\newlabel{categories-section-posets-limits@cref}{{[originalreference][0][]categories-section-posets-limits}{[1][0][]source}{}{}{}}
\makeatother
\title{A finite type module on a quasi-compact space which is a colimit of a directed system is reached at some stage}
\author{The Stacks Project Authors}
\date{stacks-project:modules.tex:local-snapshot-20261012; GNU FDL 1.2 or later}
\begin{document}
\maketitle
\noindent\textit{Editorial note.} All delivered mathematical text is the author's own, copied from the pinned Stacks Project chapter source sources/stacks-modules/pinned-original-source.tex; no statement or proof step was rewritten, reordered or removed.\par\medskip
\noindent\textit{Editorial note.} Preamble provenance: the Stacks Project keeps one single shared preamble for every chapter; this card takes it directly from the staged sources/stacks-modules/preamble.tex as three verbatim spans (lines 8-12, 14-18 and 154-217), so no derived file was written and nothing was added or redefined. The wrapper class and the base packages are supplied by the assembler.\par\medskip
\noindent\textit{Editorial note.} Bibliography provenance: neither the counted statement nor its proof contains a citation (this lane rescans the whole span before writing the card), so the card carries no bibliography and no entry was reconstructed.\par\medskip
\noindent\textit{Editorial note.} Selection provenance: this card was opened under the file-based control contract (CONTROL.md), with the mechanical pre-card gate run on its statement and proof interval before the card was written.\par\medskip
\noindent\textit{Allowed standard prerequisites (editorial).}\begin{itemize}
\item Colimits of sheaves of modules over a directed preordered set, and the description of the colimit through its sections. The author's section on modules of finite type, in which the counted statement occurs, with the general theory of such colimits in a cross-chapter section. Fixes the system whose colimit is the given module.
\item Local generation of a module of finite type by finitely many sections, and the lifting of a section of the colimit to a stage of the system. Used inside the counted proof at source lines 1159-1165. Produces finitely many sections of the colimit and indices at which they appear.
\end{itemize}\medskip
\ifdefined\hypertarget\hypertarget{sourcepointer-1}{}\fi
\noindent\textit{Source reference.} Original reference categories-section-posets-limits: sources/stacks-modules/pinned-original-source.tex, source lines 1148--1148. Navigation pointer only: the counted text refers to the Stacks item categories-section-posets-limits, which belongs to another chapter of the project and is not part of this card. The pointer carries the target's own label as its number because this chapter's staged snapshot does not contain the other chapter's tag table; its content is not reproduced here.\par\medskip
\begin{lemma}
\label{lemma-finite-type-quasi-compact-colimit}
Let $X$ be a ringed space. Let $I$ be a preordered set and
let $(\mathcal{F}_i, f_{ii'})$ be a system over $I$ consisting of sheaves
of $\mathcal{O}_X$-modules (see
Categories, Section \ref{categories-section-posets-limits}).
Let $\mathcal{F} = \colim \mathcal{F}_i$ be the colimit. Assume
(a) $I$ is directed,
(b) $\mathcal{F}$ is a finite type $\mathcal{O}_X$-module, and
(c) $X$ is quasi-compact. Then there exists an $i$ such that
$\mathcal{F}_i \to \mathcal{F}$ is surjective.
If the transition maps $f_{ii'}$ are injective
then we conclude that $\mathcal{F} = \mathcal{F}_i$ for some $i \in I$.
\end{lemma}

\begin{proof}
Let $x \in X$. There exists an open neighbourhood
$U \subset X$ of $x$ and finitely many sections
$s_j \in \mathcal{F}(U)$, $j = 1, \ldots, m$ such that
$s_1, \ldots, s_m$ generate $\mathcal{F}$ as $\mathcal{O}_U$-module.
After possibly shrinking $U$ to a smaller open neighbourhood of $x$
we may assume that each $s_j$ comes from a section of $\mathcal{F}_i$
for some $i \in I$.
Hence, since $X$ is quasi-compact we can find a finite open
covering $X = \bigcup_{j = 1, \ldots, m} U_j$, and for each $j$
an index $i_j$ and finitely many sections $s_{jl} \in \mathcal{F}_{i_j}(U_j)$
whose images generate the restriction of $\mathcal{F}$ to
$U_j$. Clearly, the lemma holds for any index $i \in I$ which
is $\geq$ all $i_j$.
\end{proof}

\end{document}
